\chapter{Experimental design}
\label{ch:design}

\keybox{why the experiment is shaped like this}{%
Two ideas govern everything in this chapter.

The first is that a comparison is only as good as what it holds fixed. If you
compare two published methods by running each author's repository, the difference
between the two numbers contains the difference between two dataloaders, two
metric implementations and two perceptual-loss backbones, and you cannot tell
which part is the method. So every configuration here is the same checkout with
different flags, and the difference between two numbers is the flags.

The second is that a prediction made after seeing the data is not a prediction.
The acceptance levels, the decision gates, the thresholds and the compute budget
were all written down before any run, and are reproduced here unchanged. Two of
those gates could have ended the project --- if the resolution limit already binds
everywhere, or if real capture geometry turns out to be effectively isotropic,
then the method has nothing to act on. Both were designed to be cheap and to be
run first, before any method code existed. Chapter~\ref{ch:results} reports what
they returned.}

This chapter states the design: what is held fixed, what is varied, what would
count as success and what would count as failure. The commands that carry it
out, and the machinery that runs them, are Appendix~\ref{app:harness}.

\section{Compute envelope and what it forbids}
Free tier: Kaggle $2\times$T4, 16\,GB each, compute capability 7.5. 3DGS states
that 24\,GB is required for paper-quality training, so numerical identity with
the published tables is unavailable for the largest scenes. Three acceptance
levels are declared in advance.

\begin{table}[htbp]
  \centering
  \caption{Acceptance levels, declared before any run.}
  \label{tab:levels}
  \small
  \begin{tabular}{llp{0.42\linewidth}l}
    \toprule
    level & claim & criterion & where attainable \\
    \midrule
    L1 & numerical reproduction & mean PSNR within $0.20$\,dB, SSIM $0.005$, LPIPS(VGG) $0.010$ of published, at the paper's resolution, iterations and split & Blender synthetic \\
    L2 & phenomenon reproduction & the claimed effect reappears with correct sign, ordering and order of magnitude under a stated deviation & everywhere \\
    L3 & harness parity & own re-run baselines are the control for all method comparisons; published numbers cited, never used as the control & everywhere; governs the thesis tables \\
    \bottomrule
  \end{tabular}
\end{table}

L3 is standard practice: 3DGS itself marks its Mip-NeRF 360 row as copied and
reports its own re-run separately ($27.58$ against the copied $27.69$). An
independent prior on implementation identity is the ${\sim}0.6$\,dB spread
between faithful splatting implementations on the same benchmark --- three times
the L1 tolerance.

\section{Reproduction, and the single-codebase decision}
\label{sec:onecodebase}

\begin{figure}[htbp]
  \centering
  \includegraphics[width=0.99\linewidth]{figures/d8-two-arms.pdf}
  \caption{The construction every comparison in this thesis rests on. Four
  configurations, one checkout, and the only thing that differs between any two
  of them is the band-limit. Running each published method from its own
  repository would have been easier and would have meant that the difference
  between two numbers also contained the difference between two dataloaders, two
  metric implementations and two perceptual backbones, with no way to say which
  part was the method. Note that the published figures are used only as targets
  to check the reproduction against; every comparison reported in
  Chapter~\ref{ch:results} is between two of these columns.}
  \label{fig:twoarms}
\end{figure}
Both baselines run inside the Mip-Splatting codebase. The 3DGS arm is that
codebase with the 3D filter zeroed and \texttt{kernel\_size}~$=0.3$, which by
Proposition~\ref{prop:reduction} is 3DGS exactly.

This is forced as well as desirable. The multi-scale benchmark data is written
in mip-NeRF's multiscale format, and only Mip-Splatting's loader can read it;
the original 3DGS repository cannot open the dataset at all --- it falls
through every branch of the dispatch and dies. It also could not work in
principle: the multi-scale loader sets the field of view \emph{per frame},
whereas the Blender reader takes one global \texttt{camera\_angle\_x}.

The benefit is methodological: both arms share one dataloader, one metrics
implementation and one LPIPS backbone (VGG), so any difference between them is
the method rather than the plumbing. The diff that separates the arms is
enumerated in the kernel and asserted on every run, and a change outside it
aborts the session before any training starts.

\keybox{sharing a codebase shares more than the plumbing, and that is the trap}{%
``The 3D filter zeroed and \texttt{kernel\_size}~$=0.3$'' makes the arm 3DGS in
its \emph{3D} band-limit, exactly, and says nothing about its \emph{2D} one.
Both arms inherited Mip-Splatting's screen-space opacity compensation, which
3DGS does not have, so for the first three rungs the arm labelled 3DGS was
anti-aliasing better than 3DGS --- and the between-arm gap, the quantity this
whole construction exists to protect, carried a residual anti-aliasing term.
Section~\ref{sec:g2} is the diagnosis and the re-measurement. The lesson is not
that the single-codebase decision was wrong, because it is forced; it is that
``both arms share everything but the diff'' is a claim about what the diff
\emph{omits}, and omissions are exactly what a diff does not show.}

\paragraph{The arms, as run.} Arm A is Mip-Splatting unmodified. Arm B is the
same codebase with \texttt{filter\_3D} identically zero, \texttt{kernel\_size}
$=0.3$, \emph{and} the 2D opacity compensation disabled in the rasteriser --- all
three, which together are vanilla 3DGS's band-limit semantics in both domains.
The third was added after \S\ref{sec:g2} found it missing.

\section{The ladder}
\begin{table}[htbp]
  \centering
  \caption{The rungs, their budgets and their exit criteria. GPU-hours for R0
  and R1 are this project's measurements; the remainder are the estimates they
  replace.}
  \label{tab:ladder}
  \small
  \begin{tabular}{llp{0.44\linewidth}}
    \toprule
    rung & run & exit criterion \\
    \midrule
    R0 & smoke --- one scene, 7\,K, both arms, plus the twelve pre-flight assertions & toolchain fixed; estimates replaced by measurement \\
    R1 & Blender single-scale train $\rightarrow$ multi-scale test, 8 scenes $\times$ 2 arms & gate G2: $\tfrac18$-scale gap exceeds 9\,dB \\
    R2 & Blender multi-scale train and test & Mip-Splatting ahead at every scale; L1 within $0.20$\,dB \\
    R3 & real scenes that fit 16\,GB: 360 indoor $\times4$, T\&T $\times2$, DB $\times2$ & per-scene table; OOMs recorded, never worked around \\
    R5 & seed spread --- 3 seeds $\times$ 2 scenes $\times$ 2 arms & $\sigma$, without which parity is not a testable claim \\
    R4 & Mip-NeRF 360 outdoor at \texttt{images\_8} --- the only authorised deviation & L2 only, separate table, resolution in the caption \\
    \bottomrule
  \end{tabular}
\end{table}

\section{The stress suite}
\label{sec:stress}

\begin{figure}[htbp]
  \centering
  \includegraphics[width=0.99\linewidth]{figures/d4-protocols.pdf}
  \caption{The six capture protocols, drawn as what they are. The black dot is
  the object; the coloured dots are the training cameras. Only the training set
  changes --- the test set is the full orbit in every case, so the model is
  always asked about viewpoints it was not fitted from, and the six differ only
  in how much the capture knew. Two of them look alike on purpose: \emph{mixed
  focal} has the same camera positions as the full orbit and differs only in
  that half the images are downsampled (drawn hollow), which is a defect of
  camera \emph{settings} rather than of geometry; and \emph{grazing} keeps a wide
  span but selects only high-incidence views, which a top-down drawing cannot
  show. The bottom row is where the capture becomes genuinely
  ill-conditioned.}
  \label{fig:protocols}
\end{figure}
Equation~\eqref{eq:cap} predicts that the method's advantage is a function of
the angular span of the training cameras. Testing that requires captures of
varying conditioning, which are constructed as camera subsets of images already
downloaded --- no new data, no extra GPU time. Each protocol is versioned as a
JSON index list and carries its \emph{measured} span, because a subset built
with a nominal target of $6^\circ$ that comes out at $20^\circ$ is not the
capture the prediction was made for, and the prediction must then be re-evaluated
at the span that was actually built rather than the one that was asked for.

\input{generated/table-protocols}

\section{Three instruments, attached to the baseline runs}
\label{sec:instruments}
All read-only with respect to the optimiser, so the reproduction remains a
reproduction. Together they test the method's premise before the method exists.

\begin{table}[htbp]
  \centering
  \caption{The instruments and what each decides.}
  \label{tab:instrumentplan}
  \small
  \begin{tabular}{lp{0.34\linewidth}p{0.42\linewidth}}
    \toprule
    & measurement & what it decides \\
    \midrule
    I-1 & distribution of \texttt{filter\_3D} at end of training, per scene, against scene scale and distance to nearest camera & whether $s^2/\lambda_i$ ever exceeds $f_k^2$. If the floor binds almost everywhere, Proposition~\ref{prop:reduction} applies universally, the method collapses to the baseline, and the project must pivot. The cheapest possible falsification, and a by-product of a run already scheduled. \\
    I-2 & $\lambda_3/\lambda_1$ per primitive on trained, unmodified checkpoints, per capture protocol; pure PyTorch, no CUDA & whether the regime table of Equation~\eqref{eq:cap} holds on real scenes. Turns \S\ref{sec:anisotropy} from prediction into evidence before any method code is written. \\
    I-3 & per-scene mean and standard deviation across seeds, both baselines & the $\sigma$ that makes ``parity'' a quantity. Parity is then reported as within $k\sigma$ of the self-run baseline. \\
    \bottomrule
  \end{tabular}
\end{table}

One detail here has a real failure mode. A primitive that no training camera
sees has no distance to any camera, and the published code gives it the largest
distance in the scene rather than leaving it undefined. Those primitives must be
excluded from the instrument's histogram, or they pile up in its right tail and
read as though the resolution floor were binding hard --- which is precisely the
wrong answer at gate G3, and an encouraging-looking one.

\section{Metrics, baselines, statistics}
Quality is PSNR, SSIM and LPIPS (VGG) per scene and averaged, at test scales
$1\times$, $\tfrac12$, $\tfrac14$, $\tfrac18$. Cost is primitive count, model
MB, peak training VRAM, wall-clock and render fps --- same GPU, same run.
Diagnostics are the fraction of primitives where the estimation term exceeds the
floor, and the mean anisotropy of $\Sigma_{\mathrm{filt}}$; this is what makes a
parity result interpretable rather than mysterious. Statistics are three seeds
minimum, mean $\pm$ standard deviation, paired per-scene deltas.

\paragraph{One measurement the published code cannot make.} The multi-scale
test set holds all four resolutions mixed together, and the evaluation that ships
with both methods averages the whole mixture into a single number. But the papers
report four numbers, one per scale, and so does this thesis --- the entire
argument is about what happens as the test resolution leaves the training
resolution, which a single pooled average hides completely.

The four groups therefore have to be recovered. They can be, exactly and without
re-rendering, because the loader preserves the order in which the test views were
written and that order records which scale each view came from. Re-rendering at
each scale instead would \emph{resample} the images and answer a different
question, so it is not done. The recovery is checked rather than trusted: the
count-weighted mean of the four recovered groups must reproduce the pooled
number the shipped evaluation reports, and that is asserted on every scene of
every run.
